\scanpage{14}
\mspage{12}
d'où un $G$ schéma en groupes sur $K$. Soit $I_G$ le ``schéma des invariants'', isomorphe si $\unsure{C|G}$ connexe \inserted{à} $T/W$ ($T$ tore maximal \unsure{?|de $G$}, $W$ son groupe de Weyl).

Alors on trouve un élément $f_s$ de $I_G(Q)$. (NB. $I_G$ est en fait défini sur $Q$). On trouve ainsi une famille d'éléments $(f_s)_{s\in S^0}$, où $S^0$ désigne l'ens des pts de $S$ tels que $k(s)$ soit fini : famille des frobenius. Compatibilité avec les hom
\[
G\deleted{\unsure{(K)}}\longrightarrow\pi_1(S,\xi)
\]
et les hom
\[
\pi_1(S,\xi)\longrightarrow G(Q_\ell).
\]
% archive sheet number: 13

